\documentclass[11pt,a4paper]{article}

% ============================================================
% Soft Spaces Phase 2 -- comprehensive research manuscript
% Generated as a long-form technical draft
% ============================================================

\usepackage[utf8]{inputenc}
\usepackage[T1]{fontenc}
\usepackage[english]{babel}
\usepackage{lmodern}
\usepackage{microtype}
\usepackage{amsmath,amssymb,amsfonts,mathtools}
\usepackage{bm}
\usepackage{physics}
\usepackage{braket}
\usepackage{booktabs}
\usepackage{longtable}
\usepackage{array}
\usepackage{multirow}
\usepackage{geometry}
\usepackage{enumitem}
\usepackage{graphicx}
\usepackage{caption}
\usepackage{subcaption}
\usepackage{xcolor}
\usepackage{hyperref}
\usepackage{authblk}
\usepackage{siunitx}
\usepackage{listings}
\usepackage{fancyhdr}
\usepackage{titlesec}
\usepackage{float}
\usepackage{placeins}

\geometry{margin=2.35cm}
\setlength{\parskip}{0.45em}
\setlength{\parindent}{0pt}
\setcounter{secnumdepth}{3}
\setcounter{tocdepth}{2}

\hypersetup{
    colorlinks=true,
    linkcolor=blue!60!black,
    citecolor=blue!60!black,
    urlcolor=blue!60!black,
    pdftitle={Soft Spaces Phase 2},
    pdfauthor={Tom Stevns}
}

\setlength{\headheight}{14pt}
\pagestyle{fancy}
\fancyhf{}
\fancyhead[L]{Soft Spaces -- Phase 2}
\fancyhead[R]{Tom Stevns}
\fancyfoot[C]{\thepage}

\lstset{
  basicstyle=\ttfamily\small,
  breaklines=true,
  frame=single,
  columns=fullflexible
}

% ------------------------------------------------------------
% Notation
% ------------------------------------------------------------
\newcommand{\Real}{\mathrm{REAL}}
\newcommand{\Null}{\mathrm{NULL}}
\newcommand{\cH}{\mathcal H}
\newcommand{\cL}{\mathcal L}
\newcommand{\cP}{\mathcal P}
\newcommand{\cQ}{\mathcal Q}
\newcommand{\cK}{\mathcal K}
\newcommand{\cS}{\mathcal S}
\newcommand{\Id}{\mathbb I}
\newcommand{\eps}{\varepsilon}
\newcommand{\Prom}{\Pi}
\newcommand{\Span}{\operatorname{span}}
\newcommand{\dist}{\operatorname{dist}}
\newcommand{\med}{\operatorname{median}}
\newcommand{\spec}{\operatorname{spec}}
\newcommand{\fro}{\mathrm F}

\title{
\textbf{Soft Spaces in Low-Qubit Hilbert Spaces -- Phase 2}\\[0.4em]
\large Dimensional Continuation, Perturbative Effective-Subspace Geometry,\\
Controlled-Lift Theorems, 12-Qubit Confirmation, and a Falsification-Driven\\
Audit of Sparse-Pauli High-Dimensional Recurrence
}

\author[1]{Tom Stevns}
\affil[1]{Independent Researcher}
\date{Revised comprehensive technical manuscript -- September 2026}

\begin{document}
\maketitle

\begin{abstract}
Phase 2 of the Soft Spaces project was designed to test whether recurrent
low-qubit structures identified in Phase 1 persist, transform, or fail when
the Hilbert-space dimension is enlarged and the numerical evidence is exposed
to progressively stronger physical and mathematical falsification tests.
The work developed from targeted 8Q-to-9Q continuation and perturbative
subspace diagnostics through 10Q, 11Q, and a predeclared checkpointed 12Q
confirmation.  The central high-Q score was formulated in terms of a
regularized Feshbach-type effective operator on an adjacent two-eigenvector
subspace,
\[
\Sigma_i^{(\eta)}
=
P_i VQ_i\,g_\eta(E_i-Q_iHQ_i)\,Q_iVP_i,
\qquad
g_\eta(x)=\frac{x}{x^2+\eta^2},
\]
followed by a normalized coherence ratio, a spectrum-matched
\(\Real\)-versus-\(\Null\) contrast, a minimum over dephasing-like and
transverse perturbation families, and a local-control prominence statistic.

The empirical Phase-2 track produced multiple positive results.  Independent
9Q perturbation tests showed recurrent advantages in spectral susceptibility
and projector retention for several frozen candidate coordinates.  Later
dimensional-transfer experiments produced a simple arithmetic continuation
rule
\[
k_{q+1}=2k_q+1,
\]
which was frozen before the final 12Q run.  The checkpointed 12Q test used
40 seeds in eight independent five-seed batches, three perturbation
repetitions, two perturbation families, fixed local controls, and a
predeclared positive-prominence gate of \(0.75\).  Both frozen 12Q predictions
passed exactly at \(6/8=0.75\).

A second part of Phase 2 then asked a stricter question: what, precisely, had
been demonstrated?  An exact non-circular ancilla-lift theorem was proved for
a controlled coupled model, showing preservation of the complete prominence
functional under a decoupled branch lift.  Perturbative stability estimates,
an ancilla-parity theorem eliminating the linear response, Schur-complement
bounds, second-order score-curvature estimates, and outward-rounded Arb/Acb
interval arithmetic produced a computer-assisted certified radius
\[
\rho_{\mathrm{cert}}=8.9\times 10^{-11}
\]
for the frozen positive controlled-model instances.  These results are
mathematically substantive, but their scope is the controlled lift, not the
independently redrawn high-Q ensemble.

The final audit materially changed the physical interpretation of the
high-Q recurrence.  A same-seed coupling audit showed that the original
dimension-dependent random-Pauli Hamiltonians are not copied ancilla lifts.
A sparse-Pauli algebra audit then showed that the fixed five-term ensemble
develops exponentially growing spectator multiplicity, while
\(k\mapsto2k+1\) exactly preserves normalized spectral rank.  Finally, a
predeclared degenerate-basis rotation test produced sign reversals in the
\(\Real-\Null\) hotspot prominence while leaving the Hamiltonians invariant
to approximately \(10^{-14}\).  Therefore the original selected ordinal
two-vector high-Q hotspots are not basis-independent observables of the
Hamiltonian and perturbation alone.

The scientific conclusion of Phase 2 is consequently twofold.  First, the
project produced a rigorous effective-subspace formalism, an exact controlled
dimensional-lift theorem, and a validated computer-assisted perturbative
certificate.  Second, the original five-term high-Q recurrence was
successfully falsified as evidence for an intrinsic transported physical
hotspot: its ordinal progression is structurally explained by normalized
spectral-rank preservation, sparse-Pauli degeneracy, and control geometry,
and the residual score is basis dependent.  This negative endpoint is not a
failure of the programme; it is the result of carrying the original claim
through a complete falsification chain.  A new physical Phase-3 model must
therefore use basis-invariant subspace observables or a sufficiently rich
Hamiltonian ensemble, matched boundary controls, and a new predeclared
holdout protocol.
\end{abstract}

\textbf{Keywords:}
Hilbert space; low-qubit quantum systems; effective Hamiltonians; Feshbach
projection; Nakajima--Zwanzig projection; random Pauli Hamiltonians;
spectrum-matched controls; spectral degeneracy; perturbation theory;
computer-assisted proof; Qiskit; falsification; basis invariance.

\tableofcontents
\newpage


% ============================================================
\section{Claims, scope, and falsification hierarchy}
% ============================================================

A central requirement of this manuscript is that different levels of
evidence are not conflated.  Phase 2 contains four logically distinct
classes of result:

\begin{enumerate}[label=\arabic*.]
\item \textbf{Observed numerical result.}
      A finite computation under a specified protocol produced a reproducible
      value or sign pattern.
\item \textbf{Exact mathematical result.}
      A theorem follows from stated assumptions without numerical
      approximation.
\item \textbf{Computer-assisted certified result.}
      A finite stored matrix instance satisfies a theorem under outward-rounded
      interval arithmetic.
\item \textbf{Falsified interpretation.}
      A stronger physical interpretation is rejected by an explicit
      counterexample or structural audit.
\end{enumerate}

These categories are intentionally kept separate throughout the paper.

\begin{table}[H]
\centering
\caption{Claim hierarchy used in Phase 2.}
\begin{tabular}{@{}p{0.21\textwidth}p{0.33\textwidth}p{0.37\textwidth}@{}}
\toprule
Category & Typical evidence & Example from this work\\
\midrule
Observed &
Frozen numerical output under specified seeds and controls &
The two frozen 12Q candidates each produced positive prominence in 6 of 8 batches.\\

Exact theorem &
Closed-form derivation under stated assumptions &
The decoupled ancilla lift preserves the effective operator, coherence ratio,
and full prominence functional.\\

Computer-assisted certificate &
Outward-rounded interval arithmetic on stored finite matrices &
The common controlled-model certified radius
\(\rho_{\mathrm{cert}}=8.9\times10^{-11}\).\\

Falsification &
Counterexample or structural mechanism incompatible with the stronger claim &
Degenerate-basis rotations reverse the prominence sign while reconstructing
the same Hamiltonian to approximately \(10^{-14}\).\\
\bottomrule
\end{tabular}
\end{table}

The strongest admissible final statement is therefore not that a universal
Soft-Space transport law has been proved.  The admissible statement is that
Phase 2 established a rigorous controlled effective-subspace theory while the
original sparse-Pauli ordinal high-Q interpretation failed a necessary
basis-invariance test.

% ============================================================
\section{Scientific aim and the role of Phase 2}
% ============================================================

The Soft Spaces programme began with an operational question rather than a
theorem: can a finite-dimensional Hilbert space contain recurrent regions
whose local structure responds differently to perturbation or open-system
evolution than a carefully matched null model?  Phase 1 established an
empirical 4Q--8Q evidence base and identified candidate coordinates, local
ridges, anchor regions, and recurrent \(\Real-\Null\) contrasts.

Phase 2 deliberately changed the standard of evidence.  The objective was no
longer merely to find more high-scoring regions.  The objectives became:

\begin{enumerate}[label=\arabic*.]
\item determine whether low-Q candidate structure continues into 9Q--12Q;
\item separate static recurrence from perturbative or dynamical protection;
\item freeze coordinates and thresholds before confirmation runs;
\item use independent seed blocks and local falsification controls;
\item translate the numerical score into operator mathematics;
\item construct an exact dimensional-lift baseline that does not encode
      hotspot labels by hand;
\item derive stability bounds away from that baseline;
\item test whether the original high-Q ensemble actually satisfies the
      hypotheses of the controlled theorem;
\item search explicitly for spectral, algebraic, indexing, and basis
      artefacts;
\item stop the original model if a physically decisive falsification is found.
\end{enumerate}

This distinction is essential.  Phase 2 should not be read as one monotonically
strengthening ``discovery''.  It is a sequence of increasingly restrictive
questions.  Some were confirmed; some were only conditionally proved; and the
strongest original physical interpretation was ultimately rejected.

% ============================================================
\section{Hilbert-space scaling and experimental objects}
% ============================================================

For \(q\) qubits,
\[
\cH_q \cong (\mathbb C^2)^{\otimes q},
\qquad
d_q=\dim\cH_q=2^q.
\]
The Phase-2 target dimensions therefore include
\[
d_8=256,\quad
d_9=512,\quad
d_{10}=1024,\quad
d_{11}=2048,\quad
d_{12}=4096.
\]

The high-Q numerical model used a sparse Hermitian random-Pauli Hamiltonian
of the form
\[
H_q=\sum_{\ell=1}^{s} a_\ell P_\ell,
\qquad s=5,
\]
where the \(P_\ell\) are \(q\)-qubit Pauli strings and the coefficients and
Pauli labels are seed generated.  The ordered spectral decomposition is
\[
H_q=\sum_{n=0}^{d_q-1}\lambda_n\ket{u_n}\bra{u_n},
\qquad
\lambda_0\le\lambda_1\le\cdots\le\lambda_{d_q-1}.
\]

A candidate adjacent spectral pair beginning at ordinal \(i\) is represented
by
\[
P_i=\ket{u_i}\bra{u_i}+\ket{u_{i+1}}\bra{u_{i+1}},
\qquad
Q_i=\Id-P_i,
\]
with pair energy
\[
E_i=\frac{\lambda_i+\lambda_{i+1}}{2}.
\]

The pair is admitted by the frozen neighbouring-gap gate when
\[
|\lambda_{i+1}-\lambda_i|<\eps_{\mathrm{neighbor}},
\qquad
\eps_{\mathrm{neighbor}}=0.05
\]
in the later high-Q confirmation protocol.

% ============================================================
\section{REAL--NULL design}
% ============================================================

The \(\Real\)-versus-\(\Null\) construction is central because it attempts to
separate spectral effects from eigenbasis structure.

For each \(\Real\) Hamiltonian,
\[
H_{\Real}=U_{\Real}\Lambda U_{\Real}^\dagger,
\]
the \(\Null\) model retains the same ordered eigenvalue multiset
\(\Lambda\) but replaces the eigenbasis by a seeded Haar-random basis,
\[
H_{\Null}=U_{\Null}\Lambda U_{\Null}^\dagger.
\]

Thus
\[
\spec(H_{\Null})=\spec(H_{\Real}),
\]
while the basis geometry is randomized.  The same perturbation family is then
applied to both models.  A positive \(\Real-\Null\) contrast is therefore
not, by construction, attributable merely to a different one-dimensional
spectrum.

This design was scientifically useful, but the later basis-invariance audit
shows an important limitation: if the \(\Real\) Hamiltonian itself has
degenerate eigenspaces and the score selects particular eigenvectors inside
those degenerate spaces, the diagonalization convention can become part of
the measured quantity.  Spectrum matching alone does not cure that issue.

% ============================================================
\section{Perturbation families and effective-subspace operator}
% ============================================================

The later frozen protocol used two physically distinct perturbation families:

\begin{itemize}
\item dephasing-like perturbations built from local \(Z\) and nearest-neighbour
      \(ZZ\) terms;
\item transverse perturbations built from local \(X\) and nearest-neighbour
      \(XX\) terms.
\end{itemize}

For a perturbation \(V\), define the regularized scalar resolvent
\[
g_\eta(x)=\frac{x}{x^2+\eta^2}.
\]
The corresponding operator-valued function is
\[
g_\eta(A)=A(A^2+\eta^2\Id)^{-1}.
\]

The second-order effective operator acting on the selected pair is
\[
\boxed{
\Sigma_i^{(\eta)}(H,V)
=
P_iVQ_i\,
g_\eta(E_i-Q_iHQ_i)\,
Q_iVP_i
}
\]
which is a regularized Feshbach-type self-energy.  Writing
\[
w_n=
\begin{pmatrix}
\bra{u_n}V\ket{u_i} &
\bra{u_n}V\ket{u_{i+1}}
\end{pmatrix},
\]
one obtains the explicit \(2\times2\) representation
\[
\Sigma_i^{(\eta)}
=
\sum_{n\notin\{i,i+1\}}
g_\eta(E_i-\lambda_n)\,w_n^\dagger w_n.
\]

This formula is one of the main mathematical objects produced by Phase 2.
It makes the relation between local spectral geometry, perturbative coupling,
and the external Hilbert-space complement explicit.


\subsection{Why the regularizer has the form \(g_\eta(x)=x/(x^2+\eta^2)\)}

The regularization is not inserted merely as a numerical convenience.  It is
the Hermitian part of a standard complex-shifted resolvent:
\[
\boxed{
g_\eta(x)
=
\frac12
\left[
\frac{1}{x-i\eta}
+
\frac{1}{x+i\eta}
\right].
}
\]
Consequently, for a self-adjoint operator \(A\),
\[
g_\eta(A)
=
\frac12
\left[
(A-i\eta\Id)^{-1}
+
(A+i\eta\Id)^{-1}
\right].
\]

This choice has three useful properties.

First, it removes the pole of the unregularized Feshbach resolvent when
\(E_i\) approaches \(\spec(Q_iHQ_i)\).  Second, it preserves Hermiticity,
because \(g_\eta(x)\) is real for real \(x\).  Third, it retains the sign of
the detuning:
\[
\operatorname{sgn} g_\eta(x)=\operatorname{sgn}x,
\]
so levels below and above the selected pair contribute with opposite signs.
This signed structure is precisely what makes the normalized coherence ratio
sensitive to constructive retention versus cancellation.

The maximal scalar magnitude is
\[
\max_x |g_\eta(x)|=\frac{1}{2\eta},
\]
attained at \(|x|=\eta\).  Therefore
\[
\|g_\eta(A)\|_{\op}\le \frac{1}{2\eta},
\]
which is also the starting point for the later perturbation bounds.

A different regularization, for example the full complex resolvent
\((x+i\eta)^{-1}\), would retain both dispersive and absorptive parts.
The present Phase-2 score deliberately uses only the Hermitian/dispersive
part because the implemented matrix functional is real-valued and designed
as a structural comparison rather than a decay-rate observable.

% ============================================================
\section{Normalized coherence ratio}
% ============================================================

The implemented score normalizes the norm of the signed sum by the sum of the
individual contribution norms:
\[
\boxed{
C_i(H,V)
=
\frac{
\left\|
\sum_{n\notin\{i,i+1\}}
g_\eta(E_i-\lambda_n)w_n^\dagger w_n
\right\|_{\fro}
}{
\sum_{n\notin\{i,i+1\}}
|g_\eta(E_i-\lambda_n)|\,\|w_n^\dagger w_n\|_{\fro}
}
}
\]
whenever the denominator is non-zero.

Because
\[
\|w_n^\dagger w_n\|_{\fro}=\|w_n\|_2^2,
\]
the triangle inequality immediately yields
\[
0\le C_i(H,V)\le1.
\]

The original implementation used terminology suggestive of
``cancellation''.  The mathematically correct interpretation is more precise:
large \(C_i\) means the signed effective contributions retain a larger coherent
resultant relative to the sum of their magnitudes; small \(C_i\) means stronger
mutual cancellation.

For one family \(f\),
\[
\Delta C_{f,i}
=
C_{f,i}^{\Real}-C_{f,i}^{\Null}.
\]
A positive value therefore means that the \(\Real\) model retains a larger
normalized effective resultant than the spectrum-matched \(\Null\) model.

\subsection{Exact invariance inside the selected pair}

For a unitary \(U\in U(2)\) acting only inside the selected pair,
\[
w_n\mapsto w_nU.
\]
Then
\[
\Sigma_i\mapsto U^\dagger\Sigma_iU,
\]
so
\[
\|\Sigma_i\|_{\fro}
\]
is unchanged, and the denominator is also unchanged because
\(\|w_nU\|_2=\|w_n\|_2\).  Hence \(C_i\) is invariant under basis rotations
\emph{within the already selected two-dimensional subspace}.

This result must not be confused with invariance under rotations in a larger
degenerate eigenspace.  The distinction becomes decisive in the final audit.

% ============================================================
\section{From branch scores to hotspot prominence}
% ============================================================

The later Phase-2 protocol used multiple robustification layers.

Let \(f\in\{Z,X\}\) label the two perturbation families, \(r\) the perturbation
repetition, \(b\) the independent batch, and \(k\) the frozen candidate
coordinate.

First,
\[
\Delta C_{f,k,b,r}
=
C^{\Real}_{f,k,b,r}
-
C^{\Null}_{f,k,b,r}.
\]
The repetition median is
\[
D_{f,k,b}
=
\med_r \Delta C_{f,k,b,r}.
\]
The robust two-family score is
\[
R_{k,b}
=
\min\{D_{Z,k,b},D_{X,k,b}\}.
\]

For the frozen local-control neighbourhood \(N_q(k)\), formed from offsets
within \(\pm10\) in steps of \(2\), the background is
\[
B_{k,b}
=
\med_{c\in N_q(k)} R_{c,b}.
\]
The batch prominence is then
\[
\boxed{
\Prom_{k,b}
=
R_{k,b}-B_{k,b}
}
\]
and the empirical positive-prominence rate is
\[
\widehat p_k
=
\frac1B
\sum_{b=1}^{B}
\mathbf 1\{\Prom_{k,b}>0\}.
\]

For the final 12Q confirmation,
\[
B=8,
\qquad
\widehat p_k\ge0.75
\]
was frozen as the decision gate.

% ============================================================
\section{Early Phase-2 continuation: 8Q to 9Q}
% ============================================================

The early Phase-2 track developed through a sequence of increasingly physical
diagnostics rather than one final score.

\subsection{Targeted continuation and local ranking}

The baseline first extended the Phase-1 candidate geometry into 9Q and later
10Q--12Q target zones.  Candidate selection was deliberately separated from
later confirmation.  Local ranking combined structural recurrence,
candidate/anchor proximity, static separation, perturbative behaviour, and
branch consistency, with the explicit intention that the scoring formula
should not be refitted after seeing the confirmation output.

\subsection{Projector-overlap tracking}

A key methodological correction was to track the perturbed two-dimensional
subspace by maximum projector overlap instead of merely choosing the closest
energy pair after perturbation.

If \(P_0\) is the original rank-two projector and \(P_\eta\) a candidate
perturbed projector, the overlap statistic can be written as
\[
\mathcal O(P_0,P_\eta)
=
\frac{1}{2}\Tr(P_0P_\eta),
\qquad
0\le\mathcal O\le1.
\]
The favourable contrast was
\[
\Delta\mathcal O
=
\mathcal O_{\Real}-\mathcal O_{\Null}>0.
\]

In parallel, a spectral-susceptibility quantity \(\chi\) was used with the
favourable convention
\[
\Delta\chi
=
\chi_{\Null}-\chi_{\Real}>0,
\]
meaning that the \(\Real\) pair is less susceptible to the perturbation than
its matched \(\Null\) counterpart.

\subsection{Independent 9Q confirmation}

One frozen 9Q confirmation used two independent seed blocks of 25 seeds each,
five perturbation strengths,
\[
\eta\in\{10^{-4},3\times10^{-4},10^{-3},3\times10^{-3},10^{-2}\},
\]
three perturbation repetitions per seed, and five perturbation terms.

The tested source coordinates included
\[
192,\;0,\;157,\;160,\;165,\;67,\;251.
\]
The predeclared descriptive threshold classified a candidate as STRONG when
both projector-overlap and susceptibility sign recurrence reached at least
\(0.80\).

Four candidates were STRONG:
\[
192,\quad165,\quad160,\quad157.
\]
For example, candidate \(192\) gave overlap-positive and
susceptibility-positive rates of \(1.0\) and \(1.0\), respectively, with
median
\[
\Delta\mathcal O\approx+0.0028,
\qquad
\Delta\chi\approx+0.3273.
\]
Candidate \(157\) gave
\[
\text{Ov+}=0.8,
\qquad
\chi+=1.0,
\]
with median
\[
\Delta\mathcal O\approx+0.0020,
\qquad
\Delta\chi\approx+0.2764.
\]

These results established a real empirical perturbative signal on independent
seeds.  They did not yet prove that the signal identified a unique intrinsic
Hilbert-space object.

% ============================================================
\section{Falsification controls, anisotropy, and local geometry}
% ============================================================

The next sequence of versions tested whether the positive perturbative effect
was candidate-specific or merely common in nearby spectral regions.

The workflow included:

\begin{itemize}
\item matched local controls;
\item local \(\Real\)-profile measurements;
\item effective-subspace perturbation analysis;
\item target confirmation on new seed blocks;
\item multi-target generalization;
\item anisotropic-direction analysis;
\item high-confidence replication;
\item local matrix-geometry diagnostics;
\item confirmation of selected target geometry;
\item path organization across the candidate neighbourhood.
\end{itemize}

The methodological principle was explicit: a candidate was scientifically
stronger only if its favourable sign recurrence exceeded the behaviour of
predeclared matched controls.  A control performing equally well was to weaken
the Soft-Space interpretation, not to be hidden.

The anisotropy work was especially important because it rejected the idea that
only a scalar perturbation ``strength'' matters.  The effective operator
contains
\[
PVQ\,g_\eta(E-QHQ)\,QVP,
\]
so the orientation of \(V\) relative to the selected subspace and its
complement is structurally relevant.  Distinct perturbation families can
therefore probe different matrix geometry even at comparable operator norm.

% ============================================================
\section{REAL--NULL cancellation, model transfer, and hotspot prediction}
% ============================================================

Later Phase-2 versions progressively simplified and stress-tested the
mechanism.

The sequence included:

\begin{enumerate}[label=\arabic*.]
\item confirmation of a \(\Real-\Null\) difference in the matrix-level
      effective contribution;
\item dimensional transfer of the same cancellation/coherence law;
\item replacement of the original perturbation generator by a local/2-local
      chain model;
\item transfer across a second perturbation family;
\item a physics-based hotspot predictor;
\item frozen hotspot confirmation;
\item dimensional hotspot transfer;
\item blind 10Q hotspot discovery;
\item 10Q confirmation;
\item an 11Q dimensional-rule test;
\item upward-validation tests over lower dimensions;
\item harmonized 7Q--9Q validation.
\end{enumerate}

The important methodological development was the move from broad descriptive
features to a small, frozen perturbative mechanism.  By the time of the
harmonized validation, the primary statistic was no longer a large composite
feature score.  It was the branch-based matrix coherence contrast and local
prominence defined above.

% ============================================================
\section{The dimensional rule and frozen geometry}
% ============================================================

The arithmetic continuation rule that emerged and was later frozen is
\[
\boxed{k_{q+1}=2k_q+1.}
\]

For the harmonized lower-dimensional tests, representative branch coordinates
were
\[
\begin{array}{c|c|c}
\text{source}\to\text{target} & \text{branch A} & \text{branch B}\\
\hline
6Q\to7Q & 15 & 47\\
7Q\to8Q & 31 & 95\\
8Q\to9Q & 63 & 191
\end{array}
\]
with the two target-space lifts of each source coordinate implemented by
\[
i=k,\qquad i=k+2^{q_{\mathrm{source}}}.
\]

It is essential to distinguish three maps:

\begin{enumerate}[label=(\alph*)]
\item the arithmetic source-coordinate map \(k\mapsto2k+1\);
\item the branch injections \(i\mapsto i\) and
      \(i\mapsto i+2^q\);
\item a physical isometry \(J:\cH_q\to\cH_{q+1}\).
\end{enumerate}

Only the first two are automatically defined by indices.  A physical
intertwiner requires operator relations that do not follow from integer
arithmetic.


% ============================================================
\section{Predeclared 12Q confirmation protocol}
% ============================================================

The final 12Q run was treated as a frozen confirmation experiment rather than
an exploratory search.  The scientific value of the result depends on
separating decisions made before execution from observations made afterward.

The following items were fixed before the confirmation run:

\begin{table}[H]
\centering
\caption{Frozen 12Q protocol.}
\begin{tabular}{@{}p{0.29\textwidth}p{0.61\textwidth}@{}}
\toprule
Protocol element & Frozen value\\
\midrule
Source coordinates &
\(k=511\) and \(k=1535\).\\

Target-space lifts &
\((511,512)\), \((2559,2560)\), \((1535,1536)\), and
\((3583,3584)\).\\

Dimensional rule &
\(k_{q+1}=2k_q+1\).\\

Local controls &
Source coordinate offsets within \(\pm10\), step \(2\).\\

Hamiltonian ensemble &
Five-term random-Pauli Hamiltonian used in the high-Q protocol.\\

Perturbation families &
Dephasing-like \(Z/ZZ\) and transverse \(X/XX\).\\

Perturbation repetitions &
Three repetitions per branch/seed/family instance.\\

Batching &
Eight independent batches of five seeds each, 40 seeds total.\\

Neighbour-gap eligibility &
\(|\lambda_{i+1}-\lambda_i|<0.05\).\\

Branch aggregation &
Median across repetitions within each perturbation family.\\

Two-family aggregation &
Minimum of dephasing and transverse \(\Real-\Null\) contrasts.\\

Local background &
Median robust score over the frozen local controls.\\

Primary statistic &
Prominence
\(\Prom_{k,b}=R_{k,b}-B_{k,b}\).\\

Decision rule &
Positive-prominence rate
\(\widehat p_k\ge0.75\), i.e.\ at least 6 of 8 batches.\\
\bottomrule
\end{tabular}
\end{table}

The key point is methodological: the candidate coordinates, control geometry,
aggregation rule, and threshold were not selected after inspection of the
12Q batch signs.  The later sparse-Pauli and basis-invariance audits were
separate falsification experiments and did not retroactively alter the
frozen 12Q decision rule.

% ============================================================
\section{Checkpointed 12Q confirmation}
% ============================================================

The final predeclared high-Q confirmation used the source coordinates
\[
k\in\{511,1535\}.
\]
The lifted target pairs were
\[
\begin{aligned}
k=511:\quad
&(511,512),\;(2559,2560),\\
k=1535:\quad
&(1535,1536),\;(3583,3584).
\end{aligned}
\]

The run used:

\begin{itemize}
\item eight independent batches;
\item five seeds per batch, hence 40 seeds;
\item three perturbation repetitions;
\item dephasing and transverse families;
\item fixed local radius \(\pm10\);
\item control step \(2\);
\item 240 checkpoint records;
\item predeclared decision gate
      \(\widehat p_k\ge0.75\).
\end{itemize}

The result was:

\begin{table}[H]
\centering
\caption{Frozen 12Q confirmation summary.}
\begin{tabular}{@{}rrrrrrr@{}}
\toprule
\(k\) & positive batches & rate & med.\ robust & med.\ background &
med.\ prominence & status\\
\midrule
511  & 6/8 & 0.750 & 0.073706 & 0.021030 & 0.038702 & confirmed\\
1535 & 6/8 & 0.750 & 0.075264 & 0.029998 & 0.043595 & confirmed\\
\bottomrule
\end{tabular}
\end{table}

The dephasing-positive rates were \(1.000\) for both candidates.  The
transverse-positive rates were \(0.875\) and \(0.750\), respectively.

At this point the empirical statement ``2/2 frozen 12Q predictions pass the
predeclared gate'' was correct.  The remainder of Phase 2 was devoted to
deciding how much physical meaning that statement could bear.

% ============================================================
\section{Qiskit translation and a gate-level witness}
% ============================================================

A Qiskit layer was added to demonstrate a route from the exact matrix
calculation toward circuit-level measurement.

For eigenstates \(\ket{p},\ket{q}\), a perturbative transition amplitude is
\[
B_{qp}
=
\bra{q}\Delta H\ket{p}.
\]
The design cross-checked selected complex amplitudes obtained by Qiskit
statevector methods against the classical matrix calculation and constructed
Hadamard-test circuits for the real and imaginary parts of representative
Pauli matrix elements.

A reduced four-qubit hardware-oriented pilot was then evaluated with
4096 shots per circuit.  It was explicitly not presented as the full 12Q
algorithm.  In ideal simulation, the normalized \(\Real-\Null\) deltas were
approximately
\[
+0.037231
\quad\text{(dephasing)},
\qquad
+0.048340
\quad\text{(transverse)},
\]
and both passed the frozen positive-delta gate.  Under the selected noisy
simulation they remained positive,
\[
+0.031144,\qquad +0.042338.
\]

This result supports the measurability of a reduced witness.  It does not
convert the complete high-Q classical statistic into a presently economical
hardware experiment.

% ============================================================
\section{Projection-operator bridge}
% ============================================================

The effective-subspace formalism has a natural connection to both Feshbach
projection and open-system projection techniques.

Given \(P\) and \(Q=\Id-P\), the unregularized Feshbach self-energy is
\[
\Sigma_P(E)
=
PVQ(E-QHQ)^{-1}QVP.
\]
The Phase-2 quantity replaces the singular resolvent by the regularized
function \(g_\eta\).

On operator space, define
\[
\cP(A)=PAP,
\qquad
\cQ=\mathcal I-\cP.
\]
Then
\[
\cP^2=\cP,
\qquad
\cQ^2=\cQ,
\qquad
\cP\cQ=\cQ\cP=0.
\]

For a Liouvillian \(\cL(t)\), the exact Nakajima--Zwanzig memory kernel has the
structural form
\[
\cK(t,s)
=
\alpha^2
\cP\cL(t)\mathcal G(t,s)\cQ\cL(s)\cP,
\]
which explicitly contains the sequence
\[
\cP\longrightarrow\cQ\longrightarrow\cP.
\]
For a time-independent generator, the Laplace transform introduces the
resolvent
\[
\int_0^\infty e^{-zt}e^{t\cQ\cL\cQ}\,dt
=
(z-\cQ\cL\cQ)^{-1},
\]
when convergent.  Thus the Hilbert-space effective operator and the
operator-space memory kernel share a genuine projection--complement--
projection architecture, while remaining mathematically different objects.

% ============================================================
\section{A deterministic prominence-stability theorem}
% ============================================================

The minimum and median operators are 1-Lipschitz in the sup norm.  This gives
a useful theorem that separates the mathematical target from the empirical
sample.

Suppose a source target \(k\) and lifted target \(k'\) have paired control
sets, and for a fixed batch
\[
|D'_{f,k'}-D_{f,k}|\le\eps_T
\]
for both perturbation families, while
\[
|D'_{f,c'}-D_{f,c}|\le\eps_C
\]
for all paired controls.  Then
\[
\boxed{
|\Prom'_{k'}-\Prom_k|
\le
\eps_T+\eps_C.
}
\]
Hence
\[
\Prom_k>\eps_T+\eps_C
\quad\Longrightarrow\quad
\Prom'_{k'}>0.
\]

If a common bound \(\eps\) applies to target and controls,
\[
\Prom_k>2\eps
\]
is sufficient.

For eight batches, if at least six source batches satisfy the corresponding
margin condition, the \(0.75\) positive-prominence gate is preserved.

This theorem identifies the real proof obligation: it is not enough to
control the target pair.  The \(\Null\) construction, both perturbation
families, and the local-control background must be controlled as well.

% ============================================================
\section{Exact controlled ancilla lift}
% ============================================================

A non-circular dimensional-lift baseline can be constructed exactly.

Let
\[
\cH_{q+1}=\mathbb C^2\otimes\cH_q,
\qquad
J_b\ket{\psi}
=
\ket{b}\otimes\ket{\psi},
\quad b\in\{0,1\}.
\]
Choose \(\Delta>0\) and define
\[
H^\uparrow
=
\ket0\bra0\otimes(H-\Delta\Id)
+
\ket1\bra1\otimes(H+\Delta\Id),
\]
with
\[
V_f^\uparrow=\Id_2\otimes V_f.
\]
If
\[
2\Delta>
\lambda_{\max}(H)-\lambda_{\min}(H),
\]
the two branches do not interleave.

For
\[
P_{i,b}^\uparrow=J_bP_iJ_b^\dagger,
\]
one obtains exactly
\[
\boxed{
\Sigma_{i,b}^{(\eta),\uparrow}
=
J_b\Sigma_i^{(\eta)}J_b^\dagger
}
\]
and therefore
\[
\boxed{
C_{i,b}(H^\uparrow,V_f^\uparrow)
=
C_i(H,V_f).
}
\]

If the \(\Real\) and \(\Null\) systems, perturbation repetitions, target, and
all controls are lifted consistently, every layer of the prominence
functional is preserved:
\[
\boxed{
\Prom_{k,b}^{\uparrow}
=
\Prom_{k,b}.
}
\]

This is an exact theorem.  It does not insert \(511\), \(1535\), or any other
hotspot label into the model; all adjacent pairs are preserved by the same
construction.  It is therefore non-circular.

Its limitation is equally exact: the original high-Q experiment did not
construct \(H_{q+1}\) by copying \(H_q\) into a decoupled ancilla branch.
The theorem proves a mechanism that \emph{can} preserve prominence, not that
the independently generated ensemble \emph{does} use that mechanism.

% ============================================================
\section{Perturbation bounds away from the exact lift}
% ============================================================

The identity
\[
g_\eta(A)
=
\frac12
\left[
(A-i\eta\Id)^{-1}
+
(A+i\eta\Id)^{-1}
\right]
\]
gives the resolvent bound
\[
\boxed{
\|g_\eta(A')-g_\eta(A)\|
\le
\frac{\|A'-A\|}{\eta^2}
}
\]
for self-adjoint \(A,A'\), together with
\[
\|g_\eta(A)\|\le\frac{1}{2\eta}.
\]

For aligned \(P,Q\), let
\[
\Sigma=PVQg_\eta(A)QVP,
\qquad
\Sigma'=PV'Qg_\eta(A')QV'P.
\]
If
\[
M=\max\{\|V\|,\|V'\|\},
\quad
\delta_V=\|V'-V\|,
\quad
\delta_A=\|A'-A\|,
\]
then
\[
\boxed{
\|\Sigma'-\Sigma\|
\le
\frac{M\delta_V}{\eta}
+
\frac{M^2\delta_A}{\eta^2}.
}
\]

For the normalized ratio \(C=N/D\), if
\[
D,D'\ge d_0>0,
\qquad
|N'-N|\le\delta_N,
\qquad
|D'-D|\le\delta_D,
\]
then
\[
\boxed{
|C'-C|
\le
\frac{\delta_N+\delta_D}{d_0}.
}
\]

Projector motion requires a separate cluster-gap control.  If \(g_*\) is the
minimum separation of the selected rank-two cluster from the remainder of the
spectrum, then a sufficient isolation condition is
\[
\|\delta H\|<\frac{g_*}{2}.
\]
With
\[
\rho=\frac{\|\delta H\|}{g_*},
\]
this is simply
\[
\rho<\frac12.
\]

% ============================================================
\section{Controlled-coupling stress test}
% ============================================================

A fixed structured perturbation of the exact ancilla lift was introduced as
\[
H(\eps)
=
H^\uparrow+\eps(X_a\otimes G),
\]
with a frozen target, frozen controls, five fixed seeds, two fixed perturbation
families, and a predeclared grid in the dimensionless coupling
\[
\rho=\frac{\|\delta H\|}{g_*}.
\]

Three of the five source seeds had positive baseline prominence; the two
negative seeds were retained.  Along the structurally safe grid
\[
\rho\le0.3<0.5,
\]
all three positive source instances satisfied the target-plus-control
prominence inequality in both branches.

Directly computed prominence remained positive even through the much larger
tested value \(\rho=10\), although beyond \(\rho=0.5\) pair identity is no
longer guaranteed by the simple isolation theorem.  This separation between
a conservative sufficient theorem and a much more stable observed path became
the motivation for the parity analysis.

% ============================================================
\section{Ancilla parity and the absence of linear response}
% ============================================================

Define
\[
\Gamma=Z_a\otimes\Id.
\]
For
\[
H(\eps)=H^\uparrow+\eps(X_a\otimes G),
\qquad
V^\uparrow=\Id_a\otimes V,
\]
one has
\[
\Gamma H(\eps)\Gamma=H(-\eps),
\qquad
\Gamma V^\uparrow\Gamma=V^\uparrow.
\]

Because the complete pair score is invariant under simultaneous unitary
conjugation of \(H\), \(V\), and the consistently tracked projector,
\[
C(\eps)=C(-\eps).
\]
The same symmetry passes through \(\Real-\Null\) subtraction, the family
minimum, the control median, and prominence:
\[
\boxed{
\Prom(\eps)=\Prom(-\eps).
}
\]
Hence, while the relevant branches remain isolated,
\[
\boxed{
\Prom(\eps)=\Prom(0)+O(\eps^2).
}
\]

This exact selection rule explains why a first-order worst-case perturbation
bound can be extremely pessimistic for the controlled path.

A numerical decomposition confirmed this structure: over the frozen
low-coupling grid, fitted response exponents ranged from approximately
\(1.999691\) to \(2.002899\), with median \(2.000099\).  The directly computed
\(+\rho\) and \(-\rho\) prominences agreed at printed numerical precision.

% ============================================================
\section{Schur-complement mechanism}
% ============================================================

Write the controlled Hamiltonian in block form,
\[
H(\eps)
=
\begin{pmatrix}
A & \eps G\\
\eps G & D
\end{pmatrix}.
\]
For an eigenvector \((\psi_A,\psi_D)^T\) with eigenvalue \(E\),
\[
\psi_D
=
-\eps(D-E)^{-1}G\psi_A,
\]
whenever the inverse exists.  Substitution gives
\[
\left[
A-E-\eps^2G(D-E)^{-1}G
\right]\psi_A=0.
\]
The induced branch self-energy is therefore
\[
\boxed{
\Sigma_A(E,\eps)
=
\eps^2G(D-E)^{-1}G.
}
\]

If \(\Gamma_{\mathrm{br}}\) is the uncoupled separation of the two branch
spectral intervals, then
\[
\dist(E,\spec(D))
\ge
\Gamma_{\mathrm{br}}-|\eps|
\]
and
\[
\boxed{
\|\Sigma_A(E,\eps)\|_2
\le
\frac{\eps^2\|G\|_2^2}{
\Gamma_{\mathrm{br}}-|\eps|
}.
}
\]
This is a continuum second-order statement and gives an operator-level reason
for the observed quadratic response.

% ============================================================
\section{Uniform analytic and curvature certificates}
% ============================================================

A first all-directions bound produced a very small common candidate coupling
radius of order
\[
\rho\sim4.24\times10^{-16}
\]
for the three positive source instances.  This was recognized as deliberately
conservative, dominated by large inverse powers of the regularizer
\(\eta=10^{-3}\), and not as an observed failure scale.

The parity-aware second-order programme then bounded first and second
derivatives of:

\begin{itemize}
\item Riesz spectral projectors;
\item regularized resolvents;
\item the effective-operator numerator;
\item the normalization denominator;
\item the full ratio \(C\);
\item target and control \(\Real-\Null\) contrasts;
\item the final prominence.
\end{itemize}

For a source prominence \(m=\Prom(0)>0\) and a uniform curvature bound
\[
|\Prom''(\eps)|\le K_{\Prom},
\]
evenness of \(\Prom\) gives
\[
\Prom(\eps)
\ge
m-\frac12K_{\Prom}\eps^2.
\]
Thus a sufficient positive-sign condition is
\[
|\eps|
<
\sqrt{\frac{2m}{K_{\Prom}}}.
\]

The double-precision continuum evaluation produced a common candidate radius
of approximately
\[
9.0\times10^{-11},
\]
subject at that stage to the need for validated outward rounding.

% ============================================================
\section{Computer-assisted Arb/Acb validation}
% ============================================================

The interval-validation stage re-evaluated the frozen constant chain with
Arb/Acb ball arithmetic at 256-bit working precision.  The validation
explicitly enclosed:

\begin{enumerate}[label=\arabic*.]
\item eigenvalues and eigenvectors of the frozen matrices;
\item rank-one and rank-two spectral gaps;
\item coupling and perturbation operator norms;
\item numerator and denominator lower bounds;
\item first- and second-derivative constants;
\item source-prominence intervals;
\item conversion from the certified \(\eps\)-radius to \(\rho\).
\end{enumerate}


\subsection{Implementation and reproducibility of the interval certificate}

The validated stage used Arb/Acb-style ball arithmetic with outward rounding
at 256-bit working precision.  The certification path was deliberately
separated from the Qiskit circuit layer: the interval proof concerns the
stored finite matrices entering the controlled effective-subspace theorem,
not shot-noise data from quantum hardware.

The validation pipeline recomputed or enclosed the quantities needed for the
deterministic bound rather than trusting double-precision point estimates.
The chain included:
\[
\text{spectral gaps}
\to
\text{projector/resolvent bounds}
\to
\text{numerator/denominator bounds}
\to
\text{score curvature}
\to
\text{prominence lower bound}
\to
\text{certified radius}.
\]

The common certified radius is defined as the minimum lower bound over the
positive frozen controlled instances:
\[
\rho_{\mathrm{cert}}
=
\min_j \underline{\rho}_j.
\]
This is why the reported common radius is
\(8.9\times10^{-11}\), even though one individual positive case admits a
larger lower bound.

For the three positive frozen source seeds, the validated results were:

\begin{table}[H]
\centering
\caption{Outward-rounded controlled-model certificates.}
\begin{tabular}{@{}rrrrr@{}}
\toprule
seed & validated interval \(\rho\) & prominence lower &
curvature upper & certified \(\rho\) lower\\
\midrule
2533001 & \(1.0\times10^{-10}\) & 0.140983066 &
\(2.19972891\times10^{21}\) & \(8.9\times10^{-11}\)\\
2533003 & \(3.0\times10^{-10}\) & 0.0736388287 &
\(5.81115634\times10^{20}\) & \(1.07914365\times10^{-10}\)\\
2533004 & \(1.0\times10^{-10}\) & 0.235286838 &
\(6.64435839\times10^{21}\) & \(8.9\times10^{-11}\)\\
\bottomrule
\end{tabular}
\end{table}

Hence
\[
\boxed{
\rho_{\mathrm{cert}}
=
8.9\times10^{-11}.
}
\]

A precision-stability audit at 192, 256, and 384 bits reproduced the same
displayed common certified lower bound.  The two negative source seeds
remained strictly negative and were retained.

The scope statement is crucial: this is a computer-assisted theorem for the
stored finite matrices of the controlled ancilla construction.  It is not a
proof for an unspecified pre-rounding random distribution, nor for the
independently redrawn 8Q--12Q ensemble.

% ============================================================
\section{Audit of the link to the independent high-Q ensemble}
% ============================================================

Having proved a rigorous controlled model, Phase 2 returned to the original
high-Q experiment and asked whether the independent random-Pauli ensemble was
actually coupled to that theorem.

A favourable same-seed 11Q--12Q audit found:

\begin{itemize}
\item the coefficient sequences were identical for all 40 same-seed
      comparisons;
\item exact copied-ancilla Hamiltonians:
      \(0/40\);
\item exactly copied Pauli terms:
      \(0/200\);
\item Hamiltonians preserving both elementary leading-ancilla computational
      sectors:
      \(0/40\);
\item every 12Q Hamiltonian contained at least one \(X\) or \(Y\) on the
      leading qubit.
\end{itemize}

Thus the direct block-ancilla intertwiner required to identify the original
ensemble with the controlled theorem is absent in this frozen audit.

The correct logical separation is therefore:
\[
\boxed{
\text{controlled ancilla theorem}
\neq
\text{original independent high-Q recurrence}.
}
\]

% ============================================================
\section{Sparse-Pauli algebra and the spectral-artifact audit}
% ============================================================

The next audit identified a deeper structural mechanism.

\subsection{Normalized spectral-rank identity}

The arithmetic rule
\[
k' = 2k+1
\]
obeys
\[
\boxed{
\frac{k'+1}{2^{q+1}}
=
\frac{k+1}{2^q}.
}
\]
Thus it exactly preserves the normalized cumulative spectral rank represented
by the ordinal \(k+1\).

The two recurrent chains lie at normalized source positions \(1/4\) and
\(3/4\).  After the two-branch lift, the four tested target pairs occur at
\[
\frac18,\quad\frac38,\quad\frac58,\quad\frac78.
\]

This identity is arithmetic, not dynamical.  Nevertheless, once spectral
multiplicity is present it can create a highly regular apparent
``continuation'' of spectral boundaries.

\subsection{Five-term Pauli algebra}

For a fixed number \(s=5\) of Pauli generators, increasing \(q\) increases the
number of spectator degrees of freedom.  The 12Q audit found:

\begin{itemize}
\item symplectic rank \(2\) with forced multiplicity \(256\): \(5/40\) seeds;
\item symplectic rank \(4\) with forced multiplicity \(512\): \(35/40\) seeds;
\item a common Pauli anticommuter, implying exact \(E\leftrightarrow -E\)
      spectral symmetry: \(40/40\);
\item regular repeat factor
      \[
      2^{12-5}=128.
      \]
\end{itemize}

The dimension dependence of the frozen local controls was especially
revealing:

\begin{table}[H]
\centering
\caption{Controls forced inside repeated spectral blocks by dimension.}
\begin{tabular}{@{}rrr@{}}
\toprule
target & regular repeat factor & controls inside repeated block / 10\\
\midrule
7Q  & 4   & 6/10\\
8Q  & 8   & 8/10\\
9Q  & 16  & 10/10\\
10Q & 32  & 10/10\\
11Q & 64  & 10/10\\
12Q & 128 & 10/10\\
\bottomrule
\end{tabular}
\end{table}

At high \(q\), the candidates lie at algebraic spectral boundaries while the
fixed \(\pm10\), step-2 controls increasingly lie inside exactly flat repeated
blocks.  Candidate and control are therefore not asymptotically of the same
spectral-boundary class.

\subsection{Dependence of the two frozen 12Q candidates}

The local-label audit found both frozen 12Q boundaries ranked first among their
11-point local candidate/control sets.  However, the two batch-prominence
vectors had correlation
\[
0.989390733,
\]
and the candidates succeeded in the same six batches and failed in the same
two batches.  The eligible seed sets were identical.

Therefore the two \(6/8\) successes must not be treated as independent
replications.  For one candidate, under a simple \(p=1/2\) sign null,
\[
\Pr[X\ge6],\quad X\sim\mathrm{Binomial}(8,1/2),
\]
equals
\[
0.14453125.
\]
That calculation is not a complete statistical model of the experiment, but
it shows why ``2/2 confirmed'' cannot legitimately be read as two independent
small-\(p\) events.

% ============================================================
\section{Degenerate-basis invariance as a necessary physical condition}
% ============================================================

Let \(E\) be a degenerate eigenvalue with eigenspace
\(\mathcal E_E\) of dimension \(m>1\).  The Hamiltonian determines the full
spectral projector
\[
P_E
=
\sum_{j=1}^{m}\ket{e_j}\bra{e_j},
\]
but it does not determine a preferred basis
\(\{\ket{e_j}\}_{j=1}^m\).

For every \(U_E\in U(m)\),
\[
\ket{e_j'}
=
\sum_{\ell=1}^m
\ket{e_\ell}(U_E)_{\ell j}
\]
gives the same Hamiltonian and the same \(P_E\).


\subsection{Necessary invariance condition}

Let \(\mathcal F(H,V)\) denote any proposed intrinsic hotspot observable.
If \(H\) has a degenerate eigenspace \(\mathcal E_E\), and \(U_E\) acts
unitarily only inside that eigenspace, then
\[
U_EHU_E^\dagger=H.
\]
Therefore a necessary condition for intrinsic basis independence is
\[
\boxed{
\mathcal F(H,V)
=
\mathcal F(U_EHU_E^\dagger,\;U_EV U_E^\dagger)
}
\]
for physically equivalent representations, or, when \(V\) is held fixed as an
independently defined physical operator and only the eigendecomposition basis
is changed,
\[
\boxed{
\mathcal F(H,V)
\text{ must not depend on the arbitrary eigenvector coordinates used inside }
\mathcal E_E.
}
\]

Failure of this condition does not mean that the numerical implementation is
incorrect.  It means that the proposed scalar is not an intrinsic observable
of the Hamiltonian-plus-perturbation pair unless an additional physical
basis-selecting observable is specified.

If a score selects only the ordinal pair
\[
\Span\{\ket{u_i},\ket{u_{i+1}}\}
\]
inside a larger degenerate eigenspace, then the selected rank-two projector is
not determined by \(H\) alone.  A physically intrinsic observable must
therefore be invariant under all such block rotations, or it must specify an
independent physical observable that selects a basis.

This is the decisive distinction between:

\begin{itemize}
\item invariance under \(U(2)\) rotations \emph{inside an already selected
      two-dimensional subspace}, which \(C_i\) has; and
\item invariance under rotations in the \emph{larger degenerate eigenspace}
      that alter which two-dimensional slice is selected, which the original
      ordinal protocol does not have in general.
\end{itemize}

% ============================================================
\section{Frozen basis-invariance falsification}
% ============================================================

The final test instantiated that issue at 8Q, where complete calculations are
tractable.

Before execution it fixed:

\begin{itemize}
\item 12 Hamiltonian seeds;
\item source coordinates \(31\) and \(95\);
\item both target branches;
\item local radius \(\pm10\), step \(2\);
\item dephasing \(Z/ZZ\) and transverse \(X/XX\) perturbations;
\item the same coherence ratio and prominence definitions;
\item eight deterministic Haar rotations inside every degenerate eigenspace.
\end{itemize}

The \(\Real\) and \(\Null\) degenerate eigenbases were rotated separately.
Hamiltonians, eigenvalues, perturbations, candidate coordinates, local
controls, gap gate, and regularizer were held fixed.

The largest Hamiltonian reconstruction discrepancies were approximately
\[
\|H'_{\Real}-H_{\Real}\|_2
=
1.434724\times10^{-14},
\]
and
\[
\|H'_{\Null}-H_{\Null}\|_2
=
9.209561\times10^{-15}.
\]

Among ten eligible seed-candidate cases:

\begin{table}[H]
\centering
\caption{Degenerate-basis falsification result.}
\begin{tabular}{@{}lrr@{}}
\toprule
rotation & sampled range crossed zero & baseline sign reversed\\
\midrule
\(\Real\) degenerate basis & 2/10 & 2/10\\
\(\Null\) degenerate basis & 4/10 & 4/10\\
\bottomrule
\end{tabular}
\end{table}

A single sign reversal is already a counterexample to basis-invariant hotspot
sign.  Multiple frozen counterexamples therefore establish that the original
ordinal two-vector prominence is not a well-defined observable of the
degenerate Hamiltonian and perturbation alone.


This should be stated in its strongest precise form:
\[
\boxed{
\text{basis dependence}
\quad\Longrightarrow\quad
\text{the original ordinal-pair score is not intrinsic to }(H,V).
}
\]
The appropriate conclusion is not that the Hamiltonian hotspot
``disappears'' under rotation.  Rather, the selected ordinal two-vector slice
is not uniquely defined by the degenerate Hamiltonian, and the prominence sign
can change when that arbitrary slice changes while the Hamiltonian itself
remains physically the same.

% ============================================================
\section{Final logical status of the original high-Q recurrence}
% ============================================================

The final Phase-2 interpretation is not that the numerical outputs were
``wrong''.  The outputs accurately record the implemented algorithm and its
diagonalization convention.

What changed is the evidential interpretation.

The high-Q coordinate recurrence is explained by the conjunction of:

\begin{enumerate}[label=\arabic*.]
\item exact normalized spectral-rank preservation under
      \(k\mapsto2k+1\);
\item exponentially growing spectator multiplicity in the fixed five-term
      Pauli algebra;
\item candidates lying on repeated-algebra boundaries while the frozen local
      controls move increasingly deep inside repeated blocks;
\item a residual \(\Real-\Null\) pair score that is not invariant under
      physically equivalent rotations in the full degenerate eigenspaces.
\end{enumerate}

Therefore the original five-term high-Q protocol does \emph{not} establish
transport of an intrinsic physical Soft-Space hotspot between independently
generated Hilbert spaces.

This conclusion is stronger, not weaker, than merely saying that a proof is
missing.  The original interpretation has an explicit mathematical mechanism
that explains the recurrence and an explicit predeclared counterexample to the
required basis invariance.

% ============================================================
\section{What Phase 2 did establish}
% ============================================================

A fair scientific summary must preserve both the positive and negative
results.

\subsection{Empirically established within the implemented protocols}

\begin{itemize}
\item Independent 9Q perturbation tests produced reproducible positive
      susceptibility and projector-overlap contrasts for several frozen
      candidates.
\item The harmonized high-Q statistic produced locally prominent recurrent
      ordinal spectral boundaries.
\item Both frozen 12Q candidate coordinates passed the predeclared
      \(6/8=0.75\) prominence gate.
\item A reduced Qiskit gate-level witness preserved positive
      \(\Real-\Null\) contrast in ideal and selected noisy simulation.
\end{itemize}

\subsection{Proved mathematically}

\begin{itemize}
\item \(0\le C_i\le1\).
\item \(C_i\) is invariant under basis rotations inside the already selected
      rank-two subspace.
\item The target/control prominence functional satisfies a deterministic
      Lipschitz stability theorem.
\item A decoupled ancilla branch lift preserves the effective operator,
      coherence ratio, and complete prominence functional exactly.
\item Regularized-resolvent and effective-operator perturbation bounds hold
      under their stated assumptions.
\item The controlled off-diagonal ancilla coupling has exact parity, forcing
      an even prominence response and eliminating the linear term.
\item The Schur complement produces an explicit quadratic branch self-energy.
\item The controlled frozen positive instances admit an outward-rounded
      computer-assisted positive-prominence radius of
      \(8.9\times10^{-11}\).
\end{itemize}

\subsection{Falsified or not established}

\begin{itemize}
\item The arithmetic rule \(k\mapsto2k+1\) is not itself a physical
      Hilbert-space intertwiner.
\item The original independently generated 11Q--12Q Hamiltonians are not
      exact copied ancilla lifts in the audited same-seed coupling.
\item The two 12Q candidate confirmations are not independent replications.
\item The fixed local controls are not boundary-matched at high dimension.
\item The original ordinal two-vector score is not invariant under rotations
      in larger degenerate eigenspaces.
\item The original five-term high-Q recurrence therefore does not establish
      an intrinsic basis-independent transported hotspot.
\end{itemize}

% ============================================================
\section{Why the negative endpoint is scientifically valuable}
% ============================================================

The strongest outcome of Phase 2 is methodological.

A weak research workflow could have stopped after the visually attractive
12Q confirmation.  Instead, the project proceeded through:

\[
\begin{aligned}
\text{empirical recurrence}
&\to \text{operator definition}
\to \text{exact baseline theorem}
\to \text{perturbation bounds}\\
&\to \text{validated certificate}
\to \text{ensemble-coupling audit}\\
&\to \text{sparse-algebra audit}
\to \text{basis-invariance falsification}.
\end{aligned}
\]

This chain separates a genuine mathematical theorem from an invalid
generalization of that theorem.  It also identifies the precise mechanism by
which a sparse high-dimensional Hamiltonian model can generate apparently
stable ordinal spectral structure without providing a basis-independent
physical subspace observable.

From a research-integrity perspective, this is preferable to increasing the
number of seeds in a protocol whose physical observable has already failed a
necessary invariance condition.

% ============================================================
\section{Requirements for a renewed physical model}
% ============================================================

A new experiment should not merely rerun the same five-term ensemble.
It must change the scientific model.

Two principled routes are available.

\subsection{Route A: complete degenerate-eigenspace projectors}

Instead of selecting an arbitrary two-vector slice, use the full spectral
projector
\[
P_E
=
\mathbf 1_{\{E\}}(H)
\]
or a physically defined cluster projector
\[
P_{\mathcal I}
=
\mathbf 1_{\mathcal I}(H)
\]
for a spectral interval \(\mathcal I\).

The effective operator becomes
\[
\Sigma_{P}^{(\eta)}
=
PVQ\,g_\eta(E_P-QHQ)\,QVP,
\]
with \(P\) now invariant under any basis rotation inside the degenerate
eigenspace.

Controls must be matched by the same algebraic class: boundary projectors
should be compared with boundary projectors, not with interior slices in flat
repeated blocks.

\subsection{Route B: physically richer Hamiltonian ensembles}

A second route is to remove the growing spectator degeneracy by allowing the
number and locality structure of Hamiltonian terms to scale with \(q\).  A
generic local model could take the form
\[
H_q
=
\sum_{j=1}^{q}
\bm h_j\cdot\bm\sigma_j
+
\sum_{\langle j,k\rangle}
\sum_{\alpha,\beta}
J^{\alpha\beta}_{jk}
\sigma_j^\alpha\sigma_k^\beta
+\cdots
\]
with a number of independent coefficients that grows with system size.

The design requirements should include:

\begin{enumerate}[label=\arabic*.]
\item predeclared ensemble family;
\item predeclared candidate observable;
\item basis invariance by construction;
\item boundary-matched control geometry;
\item independent discovery and holdout seeds;
\item effective sample size and eligibility reporting;
\item dependence-aware statistics for multiple candidate chains;
\item frozen thresholds before holdout execution.
\end{enumerate}

% ============================================================
\section{Statistical considerations}
% ============================================================

The Phase-2 decision gate was deliberately simple:
\[
\widehat p\ge0.75.
\]
For \(B=8\), this means at least six positive batches.  Such a gate is useful
as a predeclared operational rule, but it should not be confused with a full
hypothesis test.

For independent Bernoulli signs under \(p=1/2\),
\[
\Pr[X\ge6]
=
\sum_{j=6}^{8}
\binom{8}{j}2^{-8}
=
0.14453125.
\]
In the actual 12Q data, the two candidate chains were highly correlated, so
treating them as independent would exaggerate evidence.

A future protocol should therefore use either:

\begin{itemize}
\item a hierarchical model that explicitly represents candidate dependence;
\item permutation or randomization tests respecting the seed structure;
\item a predeclared family-wise error control;
\item or a primary single-candidate holdout with secondary descriptive
      candidates.
\end{itemize}

Eligibility induced by the neighbouring-gap gate must also be reported.
A batch with two eligible seeds is not statistically equivalent to one with
five independent eligible seeds.

% ============================================================
\section{Reproducibility and versioned evidence}
% ============================================================

The project was developed as a versioned computational audit trail.  The
important scientific advantage of the long v25.x sequence is not the version
number itself, but the preservation of failed and superseded experiments.

The reproducibility package should retain, at minimum:

\begin{itemize}
\item every script used for a reported result;
\item exact command-line arguments;
\item frozen seeds and batch strides;
\item output files;
\item checkpoint records;
\item initial failed coarse interval runs;
\item the final Arb/Acb validation output;
\item the same-seed coupling audit;
\item the sparse-Pauli spectral-artifact audit;
\item the degenerate-basis invariance output.
\end{itemize}

A publication repository should distinguish:

\[
\texttt{code/},\quad
\texttt{results/},\quad
\texttt{documentation/},\quad
\texttt{paper/}.
\]

The manuscript should not ask a reader to trust a verbal claim when a frozen
machine-readable output exists.

% ============================================================
\section{Suggested evidence map}
% ============================================================

\begin{longtable}{@{}p{0.18\textwidth}p{0.34\textwidth}p{0.40\textwidth}@{}}
\caption{Phase-2 evidence map.}\\
\toprule
Stage & Scientific question & Status\\
\midrule
\endfirsthead
\toprule
Stage & Scientific question & Status\\
\midrule
\endhead

v25.2--v25.5 &
Can 8Q structure guide targeted 9Q search, and can candidate subspaces be
tracked under perturbation? &
Exploratory-to-confirmatory numerical development.\\

v25.5.1--v25.7 &
Do frozen candidates retain more projector overlap or lower susceptibility
than spectrum-matched controls on independent seeds? &
Several positive candidate-specific numerical results; controls retained.\\

v25.10--v25.19 &
Does a small perturbative effective-subspace mechanism recur across targets,
families, and local geometry? &
Positive numerical evidence; matrix-level mechanism increasingly isolated.\\

v25.20--v25.30 &
Does the mechanism transfer across perturbation models and dimensions, and can
a frozen rule predict higher-Q coordinates? &
Empirical dimensional rule and harmonized frozen validation.\\

v25.31 &
Do the two frozen 12Q predictions pass the predeclared local-prominence gate? &
Yes: both 6/8 = 0.75.\\

v25.32 &
Can a reduced perturbative witness be expressed and checked at circuit level? &
Yes in ideal and selected noisy simulation; not the full 12Q calculation.\\

v25.33--v25.39 &
Can dimensional preservation be proved in a controlled non-circular model and
certified away from zero coupling? &
Yes. Exact lift theorem plus outward-rounded certified positive radius.\\

v25.40 &
Is the original independent 11Q--12Q ensemble the controlled ancilla lift? &
No in the frozen favourable same-seed audit.\\

v25.41 &
Can the coordinate recurrence be explained by the sparse Pauli algebra and
spectral-rank geometry? &
Yes for the implemented protocol: normalized rank preservation plus growing
spectator multiplicity and unmatched controls.\\

v25.42 &
Is the residual ordinal two-vector hotspot sign basis invariant inside
degenerate eigenspaces? &
No. Frozen sign-reversal counterexamples while \(H\) is unchanged to
\(\sim10^{-14}\).\\
\bottomrule
\end{longtable}


% ============================================================
\section{Limitations and non-claims}
% ============================================================

For clarity, this manuscript does \emph{not} claim any of the following:

\begin{enumerate}[label=\arabic*.]
\item that the original five-term random-Pauli ordinal recurrence is a
      universal law of finite-dimensional Hilbert spaces;
\item that the 12Q confirmation constitutes two independent replications;
\item that the controlled ancilla theorem proves the independently redrawn
      9Q--12Q ensemble;
\item that the Qiskit reduced witness is equivalent to the complete 12Q
      classical prominence computation;
\item that the computer-assisted radius is a confidence interval for a random
      distribution of Hamiltonians;
\item that an ordinal rank-two eigenspace slice inside a larger degeneracy is
      physically meaningful without an additional basis-selecting observable.
\end{enumerate}

These non-claims are not editorial qualifications added after the fact.  They
are the logical consequences of the final audit and are part of the scientific
result.

% ============================================================
\section{Discussion}
% ============================================================

Three conceptual lessons emerge.

\subsection{A recurring coordinate need not be a transported state}

The identity
\[
\frac{2k+2}{2^{q+1}}=\frac{k+1}{2^q}
\]
shows how an integer recurrence can simply preserve normalized spectral
position.  Without an operator intertwiner,
\[
H_{q+1}J=JH_q,
\]
there is no reason to interpret that recurrence as transport of an eigenstate
or eigensubspace.

\subsection{A good local score still needs a physical domain}

The normalized coherence ratio is mathematically nontrivial and useful.  Its
effective operator has a clean Feshbach structure, a bounded range, and exact
internal \(U(2)\) invariance.  None of those properties guarantees that an
ordinal pair is a physically well-defined object when embedded inside a
larger degenerate eigenspace.

This is a general lesson in computational spectral physics: invariance of a
formula with respect to coordinates inside a chosen subspace does not imply
invariance of the \emph{subspace-selection rule} itself.

\subsection{A falsification can preserve the valuable mathematics}

The basis-invariance failure does not invalidate:

\begin{itemize}
\item the projection-operator derivations;
\item the exact decoupled-ancilla theorem;
\item the parity theorem;
\item the Schur-complement mechanism;
\item the interval-arithmetic certificate.
\end{itemize}

It invalidates a proposed identification between those results and the
original independently generated sparse high-Q ensemble.

This separation is exactly why the proof ledger is useful.

% ============================================================
\section{Conclusion}
% ============================================================

Phase 2 began by asking whether Soft-Space candidates discovered in low-qubit
Hilbert spaces could be continued into 9Q--12Q and whether their apparent
stability could be given a physical and mathematical interpretation.

The empirical programme succeeded in constructing a reproducible continuation
protocol and culminated in a frozen 12Q result in which both predicted
candidate coordinates passed the predeclared \(0.75\) local-prominence gate.
A perturbative matrix mechanism was isolated, expressed as a regularized
Feshbach-type effective operator, and translated into a robust
\(\Real-\Null\) prominence statistic.

The mathematical programme then went significantly further.  It proved an
exact non-circular dimensional-lift theorem for a controlled ancilla model,
derived perturbation and prominence bounds, identified an exact parity
selection rule, reduced the controlled coupling by a Schur complement, and
produced an outward-rounded computer-assisted certificate for the frozen
positive instances.

The final falsification programme established the boundary of those positive
results.  The independently generated high-Q Hamiltonians do not satisfy the
elementary copied-ancilla intertwiner.  The fixed five-term Pauli ensemble
develops large exact degeneracies as \(q\) grows.  The arithmetic recurrence
\(k\mapsto2k+1\) preserves normalized spectral rank, and the chosen local
controls become trapped inside repeated blocks.  Most decisively, the
remaining ordinal two-vector hotspot sign can be reversed by basis rotations
inside degenerate eigenspaces while the physical Hamiltonian is unchanged.

Accordingly, the original five-term high-Q recurrence is mathematically
exhausted as evidence for an intrinsic transported physical hotspot.  The
proper Phase-2 endpoint is:

\begin{quote}
\textbf{The project has established a substantive perturbative
effective-subspace theory and a rigorous controlled dimensional-lift result,
while simultaneously falsifying the stronger interpretation that the
original sparse-Pauli ordinal high-Q recurrence is a basis-independent
physical Soft-Space transport law.}
\end{quote}

That endpoint defines the next research problem cleanly.  A new phase should
replace arbitrary ordinal slices by basis-invariant spectral projectors or use
a physically richer Hamiltonian ensemble that removes the spectator
degeneracy, together with boundary-matched controls, independent holdout data,
and a fully predeclared inference protocol.

% ============================================================
\appendix
\section{Proof-status ledger}
% ============================================================

\begin{longtable}{@{}p{0.62\textwidth}p{0.30\textwidth}@{}}
\toprule
Statement & Status\\
\midrule
\endfirsthead
\toprule
Statement & Status\\
\midrule
\endhead

\(k\mapsto2k+1\) is an exact label map &
Proved exactly.\\

The two branch injections are distinct from the arithmetic source-label map &
Proved from the frozen protocol.\\

\(0\le C_i\le1\) &
Proved exactly.\\

\(C_i\) is invariant under \(U(2)\) rotations inside the already selected
rank-two subspace &
Proved exactly.\\

Uniform target/control score bounds imply an explicit prominence bound &
Proved under stated assumptions.\\

Exact decoupled-ancilla lift preserves \(\Sigma_i\), \(C_i\), and full
prominence &
Proved exactly for the controlled lift.\\

\(\rho<1/2\) is sufficient to preserve cluster isolation under the stated
gap normalization &
Proved as a sufficient condition.\\

Controlled \(X_a\otimes G\) coupling makes prominence even in \(\eps\) &
Proved exactly for the controlled path.\\

Schur complement gives a quadratic branch self-energy &
Proved exactly under the invertibility condition.\\

Frozen controlled positive instances have common
\(\rho_{\mathrm{cert}}=8.9\times10^{-11}\) &
Computer-assisted, outward-rounded certificate for the stored matrices.\\

Original independent 11Q--12Q Hamiltonians are exact copied ancilla lifts &
Refuted for all 40 frozen same-seed audit cases.\\

High-Q ordinal recurrence can be explained by normalized-rank preservation,
sparse-Pauli multiplicity, and control geometry &
Established for the implemented five-term protocol.\\

The two frozen 12Q candidate chains are independent &
Refuted in the frozen data; batch-prominence correlation \(\approx0.9894\).\\

Ordinal two-vector hotspot sign is invariant under basis rotations in larger
degenerate eigenspaces &
Falsified by frozen finite counterexamples.\\

Original five-term high-Q protocol demonstrates an intrinsic transported
basis-independent hotspot &
Not established; the required interpretation is rejected by the final audit.\\
\bottomrule
\end{longtable}

% ============================================================
\section{Key equations in compact form}
% ============================================================

\begin{align}
P_i
&=
\ket{u_i}\bra{u_i}
+
\ket{u_{i+1}}\bra{u_{i+1}},
\\
Q_i&=\Id-P_i,
\\
E_i&=\frac{\lambda_i+\lambda_{i+1}}{2},
\\
g_\eta(x)&=\frac{x}{x^2+\eta^2},
\\
\Sigma_i^{(\eta)}
&=
P_iVQ_i\,g_\eta(E_i-Q_iHQ_i)\,Q_iVP_i,
\\
C_i
&=
\frac{\|\Sigma_i^{(\eta)}\|_{\fro}}
{\sum_{n\notin\{i,i+1\}}
 |g_\eta(E_i-\lambda_n)|\,\|w_n^\dagger w_n\|_{\fro}},
\\
D_{f,k,b}
&=
\med_r
\left(
C^{\Real}_{f,k,b,r}
-
C^{\Null}_{f,k,b,r}
\right),
\\
R_{k,b}
&=
\min\{D_{Z,k,b},D_{X,k,b}\},
\\
B_{k,b}
&=
\med_{c\in N_q(k)}R_{c,b},
\\
\Prom_{k,b}
&=
R_{k,b}-B_{k,b},
\\
\widehat p_k
&=
\frac1B
\sum_{b=1}^{B}\mathbf1\{\Prom_{k,b}>0\}.
\end{align}

The arithmetic identity behind the observed source-coordinate recurrence is
\[
k_{q+1}=2k_q+1
\quad\Longrightarrow\quad
\frac{k_{q+1}+1}{2^{q+1}}
=
\frac{k_q+1}{2^q}.
\]

The exact controlled lift is
\[
H^\uparrow
=
\ket0\bra0\otimes(H-\Delta\Id)
+
\ket1\bra1\otimes(H+\Delta\Id),
\qquad
V^\uparrow=\Id_2\otimes V,
\]
with
\[
\Sigma_{i,b}^{\uparrow}
=
J_b\Sigma_iJ_b^\dagger,
\qquad
C_{i,b}^{\uparrow}=C_i,
\qquad
\Prom_{i,b}^{\uparrow}=\Prom_{i,b}.
\]

For the structured coupling,
\[
H(\eps)=
\begin{pmatrix}
A&\eps G\\
\eps G&D
\end{pmatrix},
\]
and
\[
\Sigma_A(E,\eps)
=
\eps^2G(D-E)^{-1}G.
\]

% ============================================================
\section{Recommended figures for the publication version}
% ============================================================

The long-form manuscript should eventually be supplemented with the following
figures generated directly from frozen outputs:

\begin{enumerate}[label=\arabic*.]
\item Phase-2 pipeline: discovery, perturbative confirmation, dimensional
      transfer, 12Q confirmation, theorem track, falsification track.
\item Hilbert-space scaling \(8Q\to12Q\).
\item 9Q projector-overlap and susceptibility confirmation by candidate.
\item Dimensional source-coordinate chains under \(k\mapsto2k+1\).
\item 12Q target versus local-control prominence across eight batches.
\item Qiskit ideal/noisy reduced witness.
\item Controlled-ancilla prominence versus \(\rho\) with the
      \(\rho<1/2\) isolation region marked.
\item Log--log quadratic-response plot demonstrating parity exponent
      \(\approx2\).
\item Interval-certified curvature/radius table.
\item Sparse-Pauli repeated-spectrum geometry versus qubit number.
\item Local control locations relative to repeated spectral blocks.
\item Degenerate-basis rotation sign ranges for the v25.42 falsification.
\end{enumerate}

% ============================================================
\section{Reproducibility checklist}
% ============================================================

For each numerical result included in a final archival release, record:

\begin{enumerate}[label=\arabic*.]
\item script filename and version;
\item full command line;
\item Git commit hash;
\item qubit number;
\item Hamiltonian term count;
\item seed list, base seed, and batch stride;
\item number of batches and seeds per batch;
\item perturbation family and repetition count;
\item regularizer and gap threshold;
\item target coordinates and all controls;
\item output filename;
\item checkpoint filename if applicable;
\item predeclared decision rule;
\item whether the run is discovery, confirmation, theorem validation,
      diagnostic audit, or falsification.
\end{enumerate}

% ============================================================
\section{Terminology note}
% ============================================================

The term \emph{Soft Space} is used operationally in this project.  It should
not be understood as asserting a new fundamental mathematical object merely
because a numerical candidate has a high score.  In Phase 2 the term denotes
a hypothesized subspace or spectral region whose behaviour is tested against
matched controls under specified perturbations.

After the final high-Q audit, the phrase \emph{transported Soft-Space hotspot}
must not be used for the original five-term ordinal recurrence without an
explicit qualifier that the physical basis-independent interpretation was
not established and was contradicted by the degenerate-basis falsification.

% ============================================================
\begin{thebibliography}{99}

\bibitem{BreuerPetruccione}
H.-P. Breuer and F. Petruccione,
\textit{The Theory of Open Quantum Systems},
Oxford University Press, 2002.

\bibitem{Feshbach}
H. Feshbach,
A unified theory of nuclear reactions. II,
\textit{Annals of Physics} \textbf{19}, 287--313 (1962).

\bibitem{Nakajima}
S. Nakajima,
On quantum theory of transport phenomena,
\textit{Progress of Theoretical Physics} \textbf{20}, 948--959 (1958).

\bibitem{Zwanzig}
R. Zwanzig,
Ensemble method in the theory of irreversibility,
\textit{Journal of Chemical Physics} \textbf{33}, 1338--1341 (1960).

\bibitem{DavisKahan}
C. Davis and W. M. Kahan,
The rotation of eigenvectors by a perturbation. III,
\textit{SIAM Journal on Numerical Analysis} \textbf{7}, 1--46 (1970).

\bibitem{Kato}
T. Kato,
\textit{Perturbation Theory for Linear Operators},
Springer, 2nd ed., 1976.

\bibitem{NielsenChuang}
M. A. Nielsen and I. L. Chuang,
\textit{Quantum Computation and Quantum Information},
Cambridge University Press, 10th anniversary ed., 2010.

\end{thebibliography}

\end{document}
